\documentclass[11pt, a4paper]{article}
\usepackage[utf8]{inputenc}
\usepackage{geometry}
\usepackage{titlesec}
\usepackage{hyperref}
\usepackage{listings}
\usepackage{xcolor}
\usepackage{graphicx}
\usepackage{float}
\usepackage{tikz}
\usetikzlibrary{shapes.geometric, arrows}

\geometry{margin=1in}

% Code listing style
\definecolor{codegreen}{rgb}{0,0.6,0}
\definecolor{codegray}{rgb}{0.5,0.5,0.5}
\definecolor{codepurple}{rgb}{0.58,0,0.82}
\definecolor{backcolour}{rgb}{0.95,0.95,0.92}

\lstdefinestyle{mystyle}{
    backgroundcolor=\color{backcolour},   
    commentstyle=\color{codegreen},
    keywordstyle=\color{magenta},
    numberstyle=\tiny\color{codegray},
    stringstyle=\color{codepurple},
    basicstyle=\ttfamily\footnotesize,
    breakatwhitespace=false,         
    breaklines=true,                 
    captionpos=b,                    
    keepspaces=true,                 
    numbers=left,                    
    numbersep=5pt,                  
    showspaces=false,                
    showstringspaces=false,
    showtabs=false,                  
    tabsize=2
}

\lstset{style=mystyle}

\title{\textbf{System Design Document: AI-Powered Email Shield}}
\author{Development Team}
\date{\today}

\begin{document}

\maketitle
\tableofcontents
\newpage

\section{Introduction}
The \textbf{Email Shield} is a security middleware designed to detect and neutralize phishing threats in real-time. By leveraging a Random Forest Machine Learning model, the system analyzes URL structures within emails to identify zero-day phishing campaigns that traditional blacklist filters might miss. This document details the architectural design, module interactions, and function-level logic of the system.

\section{System Architecture}
The system operates on a micro-service architecture consisting of three primary layers:
\begin{enumerate}
    \item \textbf{Monitoring Layer (Gateway)}: Watches for incoming email files.
    \item \textbf{Security Layer}:
        \begin{itemize}
            \item \textbf{Heuristics}: Blocks dangerous attachments and spoofed files.
            \item \textbf{AI Predictor}: Analyzes URLs for phishing.
        \end{itemize}
    \item \textbf{Action Layer (Handler)}: Routes emails to Inbox or Quarantine based on verdicts.
    \item \textbf{Visualization Layer (Dashboard)}: Provides analytics on traffic, malware, and phishing.
\end{enumerate}

\section{Module Design & Function Interactions}

\subsection{1. Gateway Module (\texttt{email\_shield/gateway.py})}
The Gateway is the central orchestrator of the system. It continuously monitors the \texttt{incoming/} directory for new email files.

\subsubsection{Function: \texttt{scan\_file()}}
This function handles the lifecycle of a single email file. It reads the file, parses it, calls the AI predictor, and moves the file based on the result.

\begin{lstlisting}[language=Python]
def scan_file(filepath: str, predictor: PhishingPredictor, db: DBManager) -> None:
    """
    Reads a file, extracts URLs, scans them using the AI predictor,
    and moves the file to the appropriate folder.
    """
    print(f"Scanning {filepath}...")
    
    try:
        with open(filepath, 'r', encoding='utf-8') as f:
            content = f.read()
    except Exception as e:
        print(f"Error reading file: {e}")
        return

    # Extract Metadata
    urls = email_parser.extract_urls_from_text(content)
    user_email = email_parser.extract_recipient(content)
    filename = os.path.basename(filepath)
    
    if not urls:
        print(" -> No URLs found. Safe.")
        db.log_scan(filename, user_email, "SAFE", "N/A")
        move_file(filepath, INBOX_DIR)
        return

    print(f" -> Found {len(urls)} URLs. Checking against AI Model...")
    results = predictor.predict_urls(urls)
    
    is_malicious = False
    malicious_url = "N/A"
    
    for url, is_phishing in results.items():
        if is_phishing:
            print(f" [!!!] PHISHING DETECTED: {url}")
            is_malicious = True
            malicious_url = url
            break # Log the first malicious URL found
    
    if is_malicious:
        print(" -> Verdict: MALICIOUS. Moving to Quarantine.")
        db.log_scan(filename, user_email, "PHISHING", malicious_url)
        move_file(filepath, QUARANTINE_DIR)
    else:
        print(" -> Verdict: SAFE. Moving to Inbox.")
        db.log_scan(filename, user_email, "SAFE", "N/A")
        move_file(filepath, INBOX_DIR)
\end{lstlisting}

\paragraph{Interaction Flow:}
\begin{enumerate}
    \item \textbf{Input}: Receives a file path.
    \item \textbf{Parsing}: Calls \texttt{email\_parser} to get URLs.
    \item \textbf{Prediction}: Calls \texttt{predictor.predict\_urls} if URLs exist.
    \item \textbf{Action}: Logs to DB via \texttt{db.log\_scan} and moves file via \texttt{move\_file}.
\end{enumerate}

\subsection{2. Malware Heuristic Module}
Integrated into the Gateway, this module provides pre-computation security checks before the AI analysis.

\begin{lstlisting}[language=Python]
    # 1. Extension Blacklist
    DANGEROUS_EXTS = ['.exe', '.bat', '.vbs', '.scr']
    if any(att.lower().endswith(ext) for ext in DANGEROUS_EXTS):
        is_malware = True

    # 2. Magic Byte Verification (Anti-Spoofing)
    # Checks if a safe extension (.pdf) contains binary signatures
    if "Magic: 4D 5A" in content and not is_malware:
        is_malware = True
        malware_info = "Spoofed File Signature (MZ)"
\end{lstlisting}

\subsubsection{Function: \texttt{move\_file()}}
Safe file moving utility that prevents overwriting by renaming duplicates.

\begin{lstlisting}[language=Python]
def move_file(src: str, dest_folder: str) -> None:
    """
    Moves a file to the destination folder, handling name collisions.
    """
    filename = os.path.basename(src)
    dest_path = os.path.join(dest_folder, filename)
    
    # Handle duplicate names
    counter = 1
    while os.path.exists(dest_path):
        name, ext = os.path.splitext(filename)
        dest_path = os.path.join(dest_folder, f"{name}_{counter}{ext}")
        counter += 1
        
    shutil.move(src, dest_path)
\end{lstlisting}

\subsection{2. Email Parser Module (\texttt{email\_shield/email\_parser.py})}
A utility module for text extraction using Regular Expressions.

\subsubsection{Function: \texttt{extract\_urls\_from\_text()}}
Extracts all HTTP/HTTPS links from the raw email body.

\begin{lstlisting}[language=Python]
def extract_urls_from_text(text):
    """
    Finds all HTTP/HTTPS URLs in a text string.
    Returns a list of URLs.
    """
    # Regex pattern for finding URLs
    url_pattern = r'https?://(?:[-\w.]|(?:%[\da-fA-F]{2}))+[^\s]*'
    
    urls = re.findall(url_pattern, text)
    
    # Clean up trailing punctuation often caught by regex
    clean_urls = []
    for url in urls:
        if url.endswith('.') or url.endswith(','):
            url = url[:-1]
        clean_urls.append(url)
        
    return list(set(clean_urls)) # Return unique URLs
\end{lstlisting}

\subsubsection{Function: \texttt{extract\_recipient()}}
Parses the "To:" header to identify the target user.

\begin{lstlisting}[language=Python]
def extract_recipient(text):
    """
    Extracts the email address from 'To: ...' line.
    Returns 'unknown@company.com' if not found.
    """
    match = re.search(r'To:\s*([\w\.-]+@[\w\.-]+)', text)
    if match:
        return match.group(1)
    return "unknown@company.com"
\end{lstlisting}

\subsection{3. Predictor Module (\texttt{email\_shield/predictor.py})}
Wraps the Machine Learning model interactions.

\subsubsection{Method: \texttt{PhishingPredictor.predict\_urls()}}
The core inference method that converts URLs to features and queries the model.

\begin{lstlisting}[language=Python]
    def predict_urls(self, urls: List[str]) -> Dict[str, bool]:
        """
        Takes a list of URL strings and returns a malicious/safe verdict for each.
        """
        if not self.model:
            return {url: False for url in urls} # Fail open if no model

        results = {}
        for url in urls:
            # Extract features
            feat_dict = features.extract_features(url)
            # Convert to list (matching training format)
            feat_vector = [list(feat_dict.values())]
            
            # Predict (1 = Phishing, 0 = Safe)
            prediction = self.model.predict(feat_vector)[0]
            results[url] = bool(prediction == 1)
            
        return results
\end{lstlisting}

\subsection{4. Database Manager (\texttt{email\_shield/db\_manager.py})}
Handles persistence using SQLite.

\subsubsection{Method: \texttt{DBManager.log\_scan()}}
Records the result of every scan for auditing and dashboard visualization.

\begin{lstlisting}[language=Python]
    def log_scan(self, filename, user_email, verdict, url_detected="N/A"):
        """Inserts a new scan record."""
        cursor = self.conn.cursor()
        timestamp = datetime.datetime.now().strftime("%Y-%m-%d %H:%M:%S")
        cursor.execute('''
            INSERT INTO scan_logs (timestamp, filename, user_email, verdict, url_detected)
            VALUES (?, ?, ?, ?, ?)
        ''', (timestamp, filename, user_email, verdict, url_detected))
        self.conn.commit()
\end{lstlisting}

\subsection{5. Dashboard (\texttt{dashboard.py})}
A Streamlit web application for analytics. Below is the snippet for generating the "Traffic Composition" chart.

\subsubsection{Snippet: Traffic Composition Chart}
\begin{lstlisting}[language=Python]
with c1:
    st.subheader("Traffic Composition")
    verdict_counts = df['verdict'].value_counts().reset_index()
    verdict_counts.columns = ['Verdict', 'Count']
    
    chart = alt.Chart(verdict_counts).mark_arc(innerRadius=50).encode(
        theta=alt.Theta(field="Count", type="quantitative"),
        color=alt.Color('Verdict', scale=alt.Scale(domain=['SAFE', 'PHISHING'], range=['#00CC96', '#FF4B4B'])),
        tooltip=['Verdict', 'Count']
    )
    st.altair_chart(chart, use_container_width=True)
\end{lstlisting}

\subsubsection{Snippet: Scalable User Search}
Designed for organizations with 500+ employees, replacing the dropdown with a search-and-match interface.

\begin{lstlisting}[language=Python]
    # Search Input
    search_query = st.text_input("Search User by Email", placeholder="e.g. alice@company.com")
    
    # Default to debug or first available
    target_user = "debug@test.com"
    if target_user not in user_list:
        target_user = user_list[0]

    if search_query:
        # ... (Search Logic Overrides target_user) ...
        
    if target_user:
        # ... (Render Chart) ...
\end{lstlisting}

\subsubsection{Snippet: Forensics Search Logic}
Enables full-text search across all columns in the audit log.

\begin{lstlisting}[language=Python]
# Filter by Search Term
search_term = st.text_input("Search Logs", placeholder="Type to search ...")

if search_term:
    # Filter if any string column contains the search term (case-insensitive)
    filtered_df = filtered_df[
        filtered_df.apply(lambda row: row.astype(str).str.contains(search_term, case=False).any(), axis=1)
    ]
\end{lstlisting}

    ]
\end{lstlisting}

\subsubsection{Snippet: Time-Based Forensics}
Crucial for breach containment, allowing security teams to isolate traffic from specific windows (e.g., "Last Hour").

\begin{lstlisting}[language=Python]
    # Time Filtering Logic
    if time_option == "Last Hour":
        cutoff = datetime.datetime.now() - datetime.timedelta(hours=1)
        filtered_df = df[df['timestamp'] >= cutoff]
    elif time_option == "Custom Range":
        # ... (Date picker logic) ...
        filtered_df = df[(df['timestamp'] >= start_dt) & (df['timestamp'] <= end_dt)]
\end{lstlisting}

\subsubsection{Snippet: System Time Verification}
Displayed at the top of the dashboard to ensure the operator is aware of the machine's local time context during investigations.

\begin{lstlisting}[language=Python]
current_time = datetime.datetime.now().strftime("%Y-%m-%d %H:%M:%S")
st.caption(f"System Time: {current_time} (Local Machine Time)")
\end{lstlisting}

current_time = datetime.datetime.now().strftime("%Y-%m-%d %H:%M:%S")
st.caption(f"System Time: {current_time} (Local Machine Time)")
\end{lstlisting}

\subsubsection{Snippet: Forensics Table Configuration}
Displays the raw data with explicit Local Time formatting and Database IDs for reference.

\begin{lstlisting}[language=Python]
filtered_df['timestamp_str'] = filtered_df['timestamp'].dt.strftime('%Y-%m-%d %H:%M:%S')
st.dataframe(
    filtered_df[['id', 'timestamp_str', 'user_email', 'verdict', 'url_detected']],
    column_config={
        "id": st.column_config.NumberColumn("Scan ID", format="%d"),
        "timestamp_str": "Time (Local)"
    }
)
\end{lstlisting}

\section{Component Interaction Diagram}
Below describes the step-by-step interaction for a detected phishing email:

\begin{enumerate}
    \item \textbf{Gateway} detects \texttt{invoice.txt} in \texttt{incoming/}.
    \item \textbf{Gateway} reads content and calls \textbf{Email Parser} to extract URLs.
    \item \textbf{Email Parser} returns \texttt{['http://malicious.xyz']}.
    \item \textbf{Gateway} calls \textbf{Predictor} with the URL list.
    \item \textbf{Predictor} computes features and queries the Random Forest model. Returns \texttt{\{'http://malicious.xyz': True\}}.
    \item \textbf{Gateway} sees \texttt{True}, sets Verdict to \textbf{PHISHING}.
    \item \textbf{Gateway} calls \textbf{DB Manager} to log the incident.
    \item \textbf{Gateway} moves \texttt{invoice.txt} to \texttt{quarantine/}.
    \item \textbf{Dashboard} (on refresh) pulls new log from DB and updates the "Phishing Attacks" counter.
\end{enumerate}

\end{document}
